\chapter{课程导论}

本章对应课件第~1--20~页：课程信息、为什么需要计算数学、算法观念，以及学期内容安排。

\section{计算方法与优化}

课程名\emph{计算方法与优化}，学期为 2026--2027~秋季，主讲戴书洋，单位为武汉大学数学与统计学院。
这一页只给出身份信息，没有公式。
它规定后续每一页都在同一门课的记号系统下展开：先把连续模型写成有限步可执行的算法，再分析误差、稳定性和运算量。

\section{课程信息}

联系与成绩按课件给出：
邮箱 \code{shuyang\_dai@whu.edu.cn}，答疑用邮件或 QQ，QQ~群~337413555，办公室雷军科技楼~719。
总评
\[
\text{总评}=0.6\times\text{期末}+0.4\times\text{平时}.
\]
输入是百分制的期末分与平时分，输出是加权标量，取值通常在 $[0,100]$。
$0.6$ 与 $0.4$ 不是物理量，只是成绩权重；期末权重大，表示课程仍以可独立完成的书面与计算考查为主。

\section{大模型这么强，为什么还要上（数学、编程）课？}

课件用大模型对话截图提出问题：生成式模型已经能写代码、解算例，为何仍要学数学与编程。
这一页没有给出公式，答案在后续页展开：模型把问题翻译成方程之后，真正决定对错的是算法是否收敛、误差如何传播、浮点运算能否稳住。
大模型输出的是文本；本课要的是可验证的数值，取值范围由误差限与稳定性定理约束，不能用「看起来合理」代替。

\section{纯粹数学与应用数学}

\bookfig{fig-p04-math-split.png}{第4页}{数学分成纯粹数学与应用数学两条路线。}{纯粹数学：问题来自数学内部；应用数学：问题来自科学与工程，强调建模加求解}

课件把数学分成两支。
纯粹数学不以应用为目的，强调严格、抽象与结构，问题如哥德巴赫猜想、黎曼猜想；代表人举华罗庚、陈省身、丘成桐、陈景润、张益唐。
应用数学关心科学应用，路线是
\[
\text{实际对象}\to\text{数学模型}\to\text{数值问题}\to\text{算法与程序}.
\]
输入是物理、工程或数据中的对象，输出是可计算的方程与算法；中间每一步都丢掉一部分细节，因此后面必须估计误差。
问题来源举流体力学、材料科学、神经科学；代表人举冯·诺依曼、香农、维纳、纳什、华罗庚。
本课落在应用数学一侧：同一位华罗庚同时出现在两支，说明「算得出来」与「证得出来」并不互斥。

\section{数学重要性：以计算为例（计算古已有之）}

计算并非计算机发明后才出现。
天文用计算研究星球运动与历法，军事用计算分析弹道，航海用计算选择航线。
课件用第谷收集观测、开普勒从数据中找规律作为代表：
观测给出有限个点 $\{(t_i,x_i)\}$，规律是一个待定函数 $x(t)$。
输入是离散数据，输出是可外推的模型；数据有观测误差，所以后文必须讨论拟合与插值，而不是假设点落在精确曲线上。

\section{数学重要性：以计算为例（现代计算自二战开始）}

课件把现代大规模计算的起点放在第二次世界大战：研制核武器需要流体与辐射的数值模拟，破译密码需要对离散结构做穷举与统计。
这一页的图是历史场景，不进入公式。
它要说明的数量级变化是：问题规模一旦超出手算，就必须把连续方程离散化，并接受截断误差与舍入误差。

\section{数学重要性：以计算为例（计算必须算到目标精度）}

课件用课堂板书场景强调：理论公式写完之后，仍要把数值「算到该停的地方」。
对一个级数或迭代，停止准则是误差降到事先给定的容许值 $\varepsilon>0$，而不是「公式已经写出来」。
输入是算法与 $\varepsilon$，输出是满足
\[
\lvert\text{近似}-\text{真值}\rvert\le\varepsilon
\]
的有限步结果；$\varepsilon$ 的量纲与未知量相同，取值由应用精度决定。

\section{数学重要性：以计算为例（自己动手从零开始）}

课件用大规模手算的历史画面说明：没有现成黑箱时，必须从最基本的四则运算搭出算法。
后文秦九韶算法、浮点舍入、范数与条件数，都是「从零开始」的具体化：先规定运算顺序，再计算运算次数，再估计误差放大倍数。

\section{计算的变革}

电子计算机把计算从手算变成自动的有限步算术。
ENIAC 于 1946~年~2~月~14~日在美国投入使用，课件称其为世界上第一台通用计算机，占地约 $170$~平方米，重约 $30$~英吨，速度为每秒 $5000$~次加法或 $400$~次乘法。
这些数的输入是历史记载，输出是量级对照：加法比乘法大约快一个数量级，因此后文统计运算量时要把乘法次数单独列出。
计算机只改变执行速度，不取消误差：每一步仍是有限位算术。

\section{计算的用途：核武器数值模拟}

课件用核装置数值模拟作为计算的用途之一。
控制方程是连续的偏微分方程，机器只能解有限维代数问题，因此必须离散。
输入是连续场，输出是网格上的有限个未知数；网格越细，截断误差通常越小，未知数个数与运算量同时上升。
这一页的照片只说明应用场景，不提供可抄的公式。

\section{计算的用途：大飞机制造}

大飞机的气动、结构与控制同样依赖数值模拟。
与上一页相同：把偏微分方程变成大规模线性或非线性方程组。
输入是几何与工况，输出是可制造、可试验对照的数值预测；预测值的可信区间由网格、模型与舍入共同决定。

\section{计算能力的驱动力}

课件的结论是：计算能力的驱动来自需求，而不是先有机器再找问题。
国家需要核武器、大飞机、天气预报等能力时，受政治、经济与时效约束，只能靠计算在试验之前给出可检验的数量。
需求规定容许误差与时限，时限再反过来规定算法的运算量上限。

\section{基本问题}

计算方法要回答的基本问题是：
如何把数学模型写成数值问题；如何确定算法的适用范围；如何估计给定算法的精度；误差如何在计算中积累和传播；如何构造精度更高的算法；如何减少存储；如何分析运算量与运算时间。
课件把本课内容分成三块：
数值线性代数（线性方程组、特征值与特征向量），
数值逼近（插值与逼近），
优化理论与算法（如何训练神经网络）。
这三块对应后文的 $Ax=b$、用简单函数代替复杂函数、以及
\[
\min_{\theta}\,L(\theta).
\]
最后式的输入是参数 $\theta$ 与训练损失 $L$，输出是标量损失值；训练要的是使 $L$ 下降的 $\theta$，而不是把 $L$ 的表达式写得更长。

\section{学习方法}

核心是算法的原理及实现。
基本任务：理解每个算法的数学背景、原理和线索，对最基本的算法非常熟悉。
理论部分：认识算法，并能做收敛性、稳定性分析。
实现部分：各章算法的目的是实际计算，必须能上机并分析结果。
因此后文每个公式都要能回答两件事：用哪些输入算出什么输出，以及这个数在原问题里代表误差、残差、步长还是条件数。

\section{算法}

算法不是单独一条数学公式，而是由基本运算及其顺序组成的完整求解方案。
可用流程图描述；选择算法是数值计算的关键环节，不同算法有不同的优势与劣势。
若记基本运算集合为 $\mathcal{O}$、顺序为 $\prec$，则一个算法是 $(\mathcal{O},\prec)$。
输入是问题实例，输出是有限步之后的近似解；运算次数是非负整数，运行时间与它成正比，但不等于墙钟时间。

\section{算法（多项式与秦九韶）}

多项式
\[
P_n(x)=a_n x^n+a_{n-1}x^{n-1}+\cdots+a_1 x+a_0
\]
的输入是系数 $(a_0,\ldots,a_n)$ 与求值点 $x$，输出是标量 $P_n(x)$。
$P_n(x)$ 就是该多项式在 $x$ 处的函数值，量纲与 $a_0$ 相同。

方法一：先算每个 $a_i x^i$（$i=0,1,\ldots,n$）再相加。
乘法次数为
\[
1+2+\cdots+n=\frac{n(n+1)}{2},
\]
加法次数为 $n$。
左边对 $i=1,\ldots,n$ 统计「$x$ 连乘 $i$ 次」的乘法；输出是次数，单位是次，取值 $\{0,1,2,\ldots\}$。

方法二：秦九韶算法（1247）/Horner 算法（1819），
\[
P_n(x)=\bigl(\cdots((a_n x+a_{n-1})x+a_{n-2})x+\cdots+a_1\bigr)x+a_0.
\]
递归为 $b_n=a_n$，
\[
b_i=b_{i+1}x+a_i,\qquad i=n-1,n-2,\ldots,0,
\]
且 $P_n(x)=b_0$。
每步输入已有 $b_{i+1}$、同一个 $x$ 和系数 $a_i$，输出新的标量 $b_i$；
$b_i$ 是从最高次项看到第 $i$ 次的嵌套值，最后的 $b_0$ 才是多项式值。
乘法次数为 $n$，加法次数为 $n$，比方法一少一个数量级的乘法。

\begin{lstlisting}
def horner(a, x):
    # a[0] = a_0, ..., a[n] = a_n
    b = a[-1]
    for ai in reversed(a[:-1]):
        b = b * x + ai
    return b
\end{lstlisting}
对 $P(x)=2x^2+3x+4$ 即 \code{a=[4,3,2]}，\code{horner(a,5)} 得到 $69$，与 $2\cdot 25+15+4$ 相同；返回值就是 $P(5)$，不是系数。

\section{算法（Basel 问题的直接求和）}

Basel 问题为
\[
\sum_{k=1}^{\infty}\frac{1}{k^2}=\frac{\pi^2}{6}.
\]
用部分和
\[
P_n=\sum_{k=1}^{n}\frac{1}{k^2}
\]
逼近 $\pi^2/6$。
输入是截断下标 $n\in\mathbb{N}$，输出是正标量 $P_n\in(0,\pi^2/6)$；
$P_n$ 是前 $n$ 项的数值和，不是 $\pi$ 本身。

直接计算的截断误差为
\[
\frac{\pi^2}{6}-P_n=\sum_{k=n+1}^{\infty}\frac{1}{k^2}.
\]
由
\[
\frac{1}{k(k+1)}<\frac{1}{k^2}<\frac{1}{k(k-1)}\quad(k\ge n+1)
\]
得到
\[
\frac{1}{n+1}<\sum_{k=n+1}^{\infty}\frac{1}{k^2}<\frac{1}{n}.
\]
左右两端是误差的可计算上下界，量纲与 $1/k^2$ 相同，取值随 $n$ 按 $O(n^{-1})$ 下降。
因此要把误差压到 $10^{-2m}$，大致需要 $n\sim 10^{2m}$ 项，直接求和很慢。

\bookfig{fig-p17-basel-direct.png}{第17页}{直接部分和的误差（左）与对数误差（右）。}{$n$：项数；Error：$\pi^2/6-P_n$；Log Error：误差的常用对数；虚线：课件给出的 $1/n$、$1/(n+1)$ 界}

\section{算法（Basel 问题的加速求和）}

利用
\[
\sum_{k=1}^{\infty}\frac{1}{k(k+1)}=1,
\qquad
\frac{1}{k^2}-\frac{1}{k(k+1)}=\frac{1}{k^2(k+1)},
\]
改写
\[
\sum_{k=1}^{\infty}\frac{1}{k^2}
=1+\sum_{k=1}^{\infty}\frac{1}{k^2(k+1)}
\approx 1+\sum_{k=1}^{n}\frac{1}{k^2(k+1)}.
\]
新部分和的输入仍是 $n$，输出是
\[
Q_n=1+\sum_{k=1}^{n}\frac{1}{k^2(k+1)},
\]
它仍是 $\pi^2/6$ 的近似，但余项满足
\[
\sum_{k=n+1}^{\infty}\frac{1}{k^2(k+1)}
<\frac{1}{n}\sum_{k=n+1}^{\infty}\frac{1}{k(k+1)}
<\frac{1}{n^2}.
\]
$1/n^2$ 是余项上界，比直接法的 $1/n$ 高一阶。
同一 $n$ 下 $Q_n$ 比 $P_n$ 更接近 $\pi^2/6$；对数误差图上加速后的点下降明显更快。

\bookfig{fig-p18-basel-accel.png}{第18页}{加速后的误差（左）与对数误差（右），绿色点为加速方案。}{$n$：项数；蓝点：直接法；绿点：加 $1/(k^2(k+1))$ 后的余项}

\section{算法（算法优劣的影响）}

算法优劣直接影响速度和效率。
对小型问题，运算次数与内存似乎无所谓；对大型问题，它们起决定作用。
不合适的算法还会因为近似性和误差的传播、积累而损害精度，甚至使计算失败。
结合前两页：同样逼近 $\pi^2/6$，把余项从 $O(n^{-1})$ 改到 $O(n^{-2})$ 并不改问题，只改算法；
输出的误差数量级可以差几个 $10$ 的幂。

\section{上课内容}

学期按周安排为：
第~1~周基础；
第~2--4~周线性方程的直接法、迭代法、梯度法（最速下降与共轭梯度）；
第~5--6~周最小二乘；
第~7--8~周非线性方程与 SVD；
第~9--10~周插值逼近；
第~11--15~周优化。
这一页是后续各讲的课表，不是第~1~章内部小节名。
当前笔记只把「第~1~周基础」写完；其余周次等对应课件后再按页追加。
